\section{Các kịch bản demo}

Các kịch bản dưới đây sẽ được trình bày cụ thể trong phần demo. Mục tiêu là cho thấy hệ thống có thể thay đổi hành vi ứng dụng bằng cấu hình Feature Flag, đồng thời vẫn giữ được tính xác định, khả năng rollback nhanh, audit và quan sát vận hành.

\subsection{Quan sát project, flag và ứng dụng demo}

Kịch bản mở đầu giới thiệu project \texttt{demo-project}, các feature flag đã seed, các nhóm flag như \texttt{checkout-experience} và \texttt{recommendations}, sau đó mở Demo Application để cho thấy ứng dụng đang nhận quyết định từ Evaluation API thông qua SDK.

\subsection{Bật và tắt tính năng toàn cục}

Kịch bản này thay đổi cấu hình serving mode hoặc trạng thái runtime của một flag đang bật toàn cục. Demo Application được evaluate lại với cùng user context để quan sát tính năng xuất hiện hoặc bị tắt mà không cần build/deploy lại mã nguồn.

\subsection{Targeting theo role}

Kịch bản này chuyển giữa các user mẫu như beta customer, regular customer và shop admin. Demo sẽ cho thấy cùng một flag có thể bật cho nhóm người dùng có role phù hợp và tắt cho nhóm còn lại, đồng thời hiển thị \texttt{projectKey}, \texttt{flagKey}, \texttt{enabled} và \texttt{reason}.

\subsection{Rollout theo phần trăm}

Kịch bản này dùng nhiều tài khoản mẫu có \texttt{targetingKey} ổn định để minh họa percentage rollout. Một số tài khoản nhận tính năng mới, số khác vẫn dùng trải nghiệm cũ; khi evaluate lại cùng context, kết quả không thay đổi do stable hashing.

\subsection{Kill switch cấp flag và group}

Kịch bản này bật kill switch ở cấp một flag hoặc cấp group, ví dụ nhóm \texttt{checkout-experience}. Demo sẽ nhấn mạnh rằng trạng thái runtime chuyển sang Off ngay lập tức, trong khi lifecycle/configuration của từng flag vẫn được giữ lại để có thể khôi phục nhanh.

\subsection{Audit log, history và statistics}

Kịch bản này thực hiện một số mutation cấu hình và nhiều lần evaluation từ Demo Application, sau đó mở Dashboard để xem audit log, configuration history và statistics. Phần demo sẽ cho thấy actor, action, request ID, before/after snapshot và số liệu aggregate theo flag/reason.

\subsection{Fallback của JavaScript SDK}

Kịch bản này mô phỏng Evaluation API lỗi hoặc timeout. SDK \texttt{@ffp/js-sdk} sẽ trả kết quả fail-closed với \texttt{enabled=false}, giúp Demo Application tiếp tục hoạt động ở trạng thái an toàn và phân biệt được lỗi client-local với quyết định từ backend.

\subsection{Phân quyền Admin, Developer và Viewer}

Kịch bản cuối cùng chuyển giữa ba vai trò demo trên Admin Dashboard. Viewer chỉ xem dữ liệu, Developer được thao tác trong phạm vi cấu hình được phép, còn Admin có quyền quản trị đầy đủ; các mutation bị từ chối sẽ được backend bảo vệ thay vì chỉ ẩn ở giao diện.
